% =============================================================================
%  Chapter 5 — PETSc 3.24.5 Patch: section.c Indexing Bug
% =============================================================================

\chapter{PETSc Patch: \texttt{section.c} Indexing Bug}
\label{ch:petsc}

This chapter documents a bug in PETSc~3.24.5~\cite{PETSc2024} that affects
higher-order finite elements whose DMPlex point numbering does not start at
zero.  The bug was discovered during the development of simplex element support
in \texttt{fall\_n} and confirmed to be an upstream issue.


% ─────────────────────────────────────────────────────────────────────────
\section{Context: DMPlex and PetscSection}
\label{sec:dmplex_context}

PETSc's \texttt{DMPlex} represents an unstructured mesh as a directed acyclic
graph (DAG) of \emph{points}: vertices, edges, faces, and cells, numbered from
$0$ to $N-1$ (though the numbering of each \emph{stratum} may start at a
non-zero offset).

A \texttt{PetscSection} associates degrees of freedom (DOFs) with mesh points.
It uses an \emph{atlas} (\texttt{s->atlasDof[]}) indexed relative to the
chart $[\texttt{pStart}, \texttt{pEnd})$, and a \emph{constraint} mechanism
to impose Dirichlet boundary conditions by specifying which DOFs at a point
are constrained.

The function \texttt{PetscSectionSetConstraintIndices()} sets the constrained
DOF indices at a given point.  Its implementation reads:
\begin{lstlisting}[language=C]
PetscCall(PetscSectionCheckConstraints_Private(s));
/* Bug: should subtract pStart to get the atlas offset */
if (s->atlasDof[point]) {              // <-- HERE
    PetscInt cdof;
    PetscCall(PetscSectionGetConstraintDof(s, point, &cdof));
    // ...
}
\end{lstlisting}


% ─────────────────────────────────────────────────────────────────────────
\section{The Bug}
\label{sec:the_bug}

In \texttt{src/sys/utils/section.c}, at approximately line~2997, the code reads:
\begin{lstlisting}[language=C,escapechar=@]
if (s->atlasDof[@\textcolor{red}{point}@]) {
\end{lstlisting}
but the atlas array is indexed relative to \texttt{pStart}, so the correct
access is:
\begin{lstlisting}[language=C,escapechar=@]
if (s->atlasDof[@\textcolor{green!50!black}{point - s->pStart}@]) {
\end{lstlisting}

When \texttt{pStart = 0} (as happens for vertex stratum when vertices are
numbered from zero), the two expressions are identical and the bug is
invisible.  However, when a mesh's vertex stratum starts at a non-zero offset
--- which occurs in \texttt{fall\_n} because Gmsh-imported meshes place cells
before vertices in the DAG, giving the vertex stratum
$\texttt{pStart} = 8$~(for 8~hex cells) or larger --- the raw \texttt{point}
index overflows the atlas array, reading uninitialised memory.


% ─────────────────────────────────────────────────────────────────────────
\section{Manifestation in \texttt{fall\_n}}
\label{sec:manifestation}

The bug manifests during the application of Dirichlet boundary conditions
via \texttt{PetscSectionSetConstraintIndices}.  Specifically:
\begin{itemize}[nosep]
  \item For HEX27 meshes with 20 cells, the vertex stratum starts at
        \texttt{pStart} $\ge$ 20.
  \item When setting constraints on vertex DOFs, the index
        \texttt{s->atlasDof[point]} reads from memory at offset 20+
        instead of offset 0+.
  \item The read value is unpredictable: it may be zero (causing the
        constraints to be silently skipped) or non-zero (proceeding correctly
        by accident).
  \item The symptom was a segmentation fault or incorrect constraint
        application on some meshes but not others, depending on memory layout.
\end{itemize}


% ─────────────────────────────────────────────────────────────────────────
\section{The Fix}
\label{sec:the_fix}

The fix is a one-line change in
\texttt{src/sys/utils/section.c}:

\begin{lstlisting}[language=C]
// Before (buggy):
if (s->atlasDof[point]) {

// After (fixed):
if (s->atlasDof[point - s->pStart]) {
\end{lstlisting}

This was applied to the local PETSc installation at:
\begin{lstlisting}[basicstyle=\ttfamily\footnotesize,breaklines=true]
/home/sechavarriam/MyLibs/petsc-3.24.5/src/sys/utils/section.c
\end{lstlisting}

After applying the patch, PETSc was rebuilt and all \texttt{fall\_n} constraint
tests passed on both HEX27 and TET10 meshes.


% ─────────────────────────────────────────────────────────────────────────
\section{Verification}
\label{sec:petsc_verification}

The fix was verified by:
\begin{enumerate}[nosep]
  \item Running the debug output in \texttt{fall\_n}'s boundary condition
        application code, which prints:
\begin{lstlisting}[basicstyle=\ttfamily\footnotesize]
[SIEVE-DBG] section_ptr=0x... chart=[9152,23629)
[SIEVE-DBG] All 14477 section DOFs OK
[BC-DBG] Constraints count: 145
[BC-DBG] plex_id=9152 section_dof=3 constraint_size=3
  ...
\end{lstlisting}
        confirming that all 145~Dirichlet-constrained vertices are correctly
        handled with \texttt{pStart} = 9152.

  \item Comparing the Case~1 (SmallStrain + Linear Elastic) displacement field
        on HEX27 against the expected Timoshenko beam solution: the
        tip displacement matches the analytical result to within the expected
        discretisation error.

  \item Running all 7~nonlinear cases on both HEX27 and TET10 meshes without
        segfaults or constraint errors.
\end{enumerate}


% ─────────────────────────────────────────────────────────────────────────
\section{Upstream Status}

As of PETSc~3.24.5, this bug is present in the release source.  It has been
reported to the PETSc developers.  Since the bug only affects codes that:
\begin{itemize}[nosep]
  \item use \texttt{PetscSectionSetConstraintIndices} directly (rather than
        through higher-level \texttt{DMPlexSetBoundaryConditions}), and
  \item have a \texttt{PetscSection} chart with \texttt{pStart} $> 0$,
\end{itemize}
it may not have been triggered by the PETSc test suite, which may use meshes
with vertex-first ordering (\texttt{pStart} = 0 for the vertex stratum).
